\section{First-order rigidity of all nearly optimal two-label programs}
\label{sec:dot-first-order-rigidity}

The exact value asymptotic in Theorem~\ref{thm:r80-sharp-modulus}
does not by itself identify every nearly optimal program.  We now
characterize them to the first nontrivial order.  This is an author
addition to the fixed scalar family of Section~\ref{sec:r80-sharp-stability},
not a classification for arbitrary quantum instruments.  The entire
six-parameter stochastic controller class and its prescribed initial
label remain unchanged.

Keep $\beta,q,r_0,B_0,D,A,F,L,k,h,c_0$ from
\eqref{eq:r80-sharp-constants} and
\eqref{eq:r80-leading-coefficients}.  Define also
\begin{equation}\label{eq:dot-rigidity-coefficients}
 S=\frac{r_0}{\beta D},\qquad
 Z=\frac{r_0^2}{D^2},\qquad
 T=\frac{r_0^2}{\beta D^2},\qquad
 E=\frac{1-r_0}{r_0}.
\end{equation}
All these constants are positive, and $S-L=r_0/D^2>0$.
For a controller $\Theta_\varepsilon$ in
\eqref{eq:r80-six-parameters}, let $J_\varepsilon$ be its worst
model cost, with value $+\infty$ if it is improper from the
prescribed initial label in either model.

\begin{theorem}[Necessary and sufficient first-order optimality]
\label{thm:dot-first-order-controller-rigidity}
Choose any controller $\Theta_\varepsilon$ for each sufficiently
small $\varepsilon>0$.  No continuity or differentiability of this
selection is required.  Then
\begin{equation}\label{eq:dot-near-optimal-scale}
 J_\varepsilon-V_2(\varepsilon)=o(\sqrt\varepsilon)
\end{equation}
if and only if its coordinates satisfy
\begin{align}
 r&=r_0+h\sqrt\varepsilon+o(\sqrt\varepsilon),
 &s&=o(\sqrt\varepsilon),\nonumber\\
 \delta&=k\sqrt\varepsilon+o(\sqrt\varepsilon),
 &\gamma&=o(\varepsilon),\nonumber\\
 1-\alpha&=o(\sqrt\varepsilon).
 &&\label{eq:dot-optimal-controller-scales}
\end{align}
The remaining coordinate $\xi$ may be any function into $[0,1]$.
Every exact optimizer satisfies these conclusions.  Under either
equivalent condition, each model separately has cost
$B_0+c_0\sqrt\varepsilon+o(\sqrt\varepsilon)$.
\end{theorem}
\begin{proof}
Write $\lambda=\sqrt\varepsilon$.  Assume
\eqref{eq:dot-near-optimal-scale}.  The exact value asymptotic gives
$J_\varepsilon=B_0+c_0\lambda+o(\lambda)$.
For a fixed controller, the baseline model-$0$ failure matrix is
entrywise no greater than its perturbed matrix.  Termwise comparison
of the nonnegative finite-prefix cost series therefore gives
$B_0\le J_0\le J_\varepsilon$, where $J_0$ is this controller's
baseline robust cost, not the optimum over baseline controllers.

Every limit point in the compact parameter cube is a baseline
optimizer: the baseline total cost is the supremum of continuous
finite-prefix costs and is lower semicontinuous.  The complete equality
set in Lemma~\ref{lem:r80-boundary-rigidity} consequently gives
\[
 r\longrightarrow r_0,\quad s,\gamma,\delta\longrightarrow0,
 \quad\alpha\longrightarrow1,
\]
with no restriction on $\xi$.  Applying
Lemma~\ref{lem:r80-boundary-coercivity} with
$0\le J_0-B_0=O(\lambda)$ gives, uniformly along the selection,
\begin{equation}\label{eq:dot-preliminary-scales}
 r-r_0=O(\lambda),\qquad s+\delta=O(\lambda),\qquad
 \gamma=O(\lambda^2).
\end{equation}

\emph{The missing update scale.}
Let $C_1$ be the actual model-$1$ cost and let $C_1^I$ denote
its cost after setting $\alpha=1,\xi=0$ for this comparison only.
Near the baseline equality set one has $r>s$.  Both row success
probabilities are bounded away from zero, so all the following
resolvents are proper and smooth, uniformly for $\xi\in[0,1]$.
Put
\[
 U=\beta(1-r),\quad W=\beta(1-s),\quad
 a=r\alpha+q(1-r)\gamma,\quad
 b=s\xi+q(1-s)\delta,\quad
 Q=UW+Ub+Wa.
\]
The exact cost and its off-diagonal derivatives are
\begin{equation}\label{eq:dot-resolvent-derivatives}
 C_1=\frac{W+a+b}{Q},\qquad
 \partial_a C_1=\frac{(U-W)(W+b)}{Q^2},\qquad
 \partial_b C_1=\frac{(W-U)a}{Q^2}.
\end{equation}
In particular, increasing $\xi$ cannot reduce the cost, whereas
$\partial_\alpha C_1$ tends uniformly to $-Z<0$ at the
baseline equality set.  In a smaller fixed neighborhood,
$-\partial_\alpha C_1\ge Z/2$, including every intermediate
$\alpha$ between its actual value and one.  Integrating that
inequality and then reducing $\xi$ to zero yields
\[
 C_1-C_1^I\ge (Z/2)(1-\alpha).
\]
The improved cost $C_1^I$ is a smooth function of
$r,s,\gamma,\delta$.  By \eqref{eq:dot-preliminary-scales},
$C_1^I=B_0+O(\lambda)$.  Since $C_1\le J_\varepsilon=B_0+O(\lambda)$,
we obtain
\begin{equation}\label{eq:dot-alpha-scale}
 1-\alpha=O(\lambda).
\end{equation}
Convergence of $\alpha$ alone would not justify its first-order
Taylor expansion; this estimate supplies the needed rate.

\emph{All scaled subsequences.}
Take an arbitrary sequence tending to zero.  The established bounds
allow extraction of a subsequence such that
\begin{align*}
 (r-r_0)/\lambda&\longrightarrow H,&
 s/\lambda&\longrightarrow u,&
 \delta/\lambda&\longrightarrow v,\\
 \gamma/\lambda^2&\longrightarrow G,&
 (1-\alpha)/\lambda&\longrightarrow a_*,&
 \xi&\longrightarrow\zeta.
\end{align*}
Here $u,v,G,a_*\ge0$ and $\zeta\in[0,1]$.  Put $t=u+v$.
We must show $t>0$ before dividing by it.

For model $0_\varepsilon$, use the four probabilities
$h_1,h_2,t_1,t_2$ in \eqref{eq:r80-perturbed-chain}, and put
$\mathcal D=h_1h_2+h_1t_2+h_2t_1$.
Finite cost and $\alpha\to1$ force $s+\delta>0$ for small
positive $\varepsilon$: otherwise the second label is absorbing
and dark, but is entered with probability at least
$\varepsilon r\alpha>0$.  Thus $\mathcal D>0$ and the exact identity
\eqref{eq:r80-trap-cost} applies.  In the fixed neighborhood,
$\mathcal D\le C(s+\delta)$, $t_1\ge\varepsilon r\alpha$,
and $h_1-h_2$ is bounded below positively.  Consequently
\[
 C_0-\frac1{h_1}
 =\frac{(h_1-h_2)t_1}{h_1\mathcal D}
 \ge\frac{c\varepsilon}{s+\delta}
\]
for constants $C,c>0$ independent of the controller.
Also $1/h_1=B_0+O(\lambda)$.  The bounded normalized robust
cost excess therefore rules out $t=0$.

Now $\mathcal D/\lambda\to r_0t$ and
$t_1/\lambda^2\to r_0+(1-r_0)G$.  Keeping the leak contribution,
rather than dropping it as in a value lower bound, gives the exact
limit
\begin{equation}\label{eq:dot-model-zero-limit}
 \frac{C_0-B_0}{\lambda}\longrightarrow
 X=-AH+\frac{1+EG}{t}.
\end{equation}
For model $1$, the rational expression in
\eqref{eq:dot-resolvent-derivatives} has a denominator bounded
away from zero in the fixed neighborhood.  At the baseline,
its derivatives, with $\xi$ held fixed, are
\[
 \partial_r C_1=F,\qquad
 \partial_s C_1=S+T\xi,\qquad
 \partial_\delta C_1=L,\qquad
 \partial_\alpha C_1=-Z.
\]
These identities follow by direct differentiation of the displayed
resolvent; they hold uniformly for $\xi\in[0,1]$.
The $\gamma$ contribution is $O(\lambda^2)$.
The bounds \eqref{eq:dot-preliminary-scales} and
\eqref{eq:dot-alpha-scale} give a uniform Taylor remainder
$o(\lambda)$.  Hence
\begin{equation}\label{eq:dot-model-one-limit}
 \frac{C_1-B_0}{\lambda}\longrightarrow
 Y=FH+Lt+(S-L)u+Za_*+T\zeta u.
\end{equation}

Since the robust normalized gap tends to $c_0$, we have
$\max(X,Y)=c_0$.  The positive weights $F/(A+F)$ and $A/(A+F)$
give
\begin{align*}
 c_0&\ge\frac{FX+AY}{A+F}\\
 &=\frac{F/t+ALt}{A+F}
   +\frac{FEG/t+A(S-L)u+AZa_*+AT\zeta u}{A+F}\\
 &\ge\frac{2\sqrt{FAL}}{A+F}=c_0.
\end{align*}
Every additional term is nonnegative.  Equality in the
arithmetic--geometric mean forces $t=k$.  Since $E,S-L,Z>0$,
equality also forces $G=u=a_*=0$, whence $v=k$.
Equality of a positive weighted average to its maximum forces
$X=Y=c_0$.  Solving $-AH+k^{-1}=FH+Lk$ gives $H=h$.
Every convergent scaled subsequence has these same limits.
Thus the full limits hold and imply
\eqref{eq:dot-optimal-controller-scales}.  This argument imposes no
limit on $\xi$.

\emph{Sufficiency and realizability.}
Conversely assume \eqref{eq:dot-optimal-controller-scales}.
For all sufficiently small $\varepsilon$, $\delta>0$,
$r$ is bounded inside $(0,1)$, and $s$ is close to zero.
In model $0_\varepsilon$ the second label can return to the first
through a $b$-failure with probability $(1-s)\delta>0$, and
the first label has success probability $(1-\varepsilon)r>0$.
Both reachable labels therefore have finite cost.  In model $1$
both row success probabilities are bounded away from zero.
The controller is consequently proper from its prescribed start
in both models.

The exact model-$0$ expansion applies with $H=h,t=k,G=0$ and
gives $(C_0-B_0)/\lambda\to-Ah+k^{-1}=c_0$.
The uniform model-$1$ expansion has $u=a_*=0,v=k$ and gives
$(C_1-B_0)/\lambda\to Fh+Lk=c_0$, regardless of $\xi$.
Therefore $J_\varepsilon=B_0+c_0\lambda+o(\lambda)$;
Theorem~\ref{thm:r80-sharp-modulus} proves
\eqref{eq:dot-near-optimal-scale} and the stated componentwise limits.
Exact optimizers are a special case because their value difference
is zero.  This completes both implications.
\end{proof}

The theorem identifies a unique leading repair mechanism in this
family.  At the square-root scale, activating success at the second
label through $s$, adding baseline leakage through $\gamma$ at the
$\varepsilon$ scale, or losing the optimal $a$-failure update through
$1-\alpha$ all incur strictly positive leading penalties.  The
$b$-failure return $\delta$ has the forced square-root coefficient
$k$, while the initial command row has the forced correction $h$.
This is neither an exact finite-$\varepsilon$ optimizer formula nor
a universal degeneration classification, and it does not close the
original general-memory squared-sign risk or the other independent
research questions recorded below.
